% Einheit 5 – Faktorisieren (binomische-formeln.md Einheit 3) – Nr. 40–48
\einheitenkopf[e5][5 Faktorisieren]{Einheit 5 von 5 · Faktorisieren}
\zweigzeile{Hier lernst du, eine Summe mit den binomischen Formeln in ein Produkt zu verwandeln · neu in diesem Jahr · keine P10-Aufgabe · baut auf: Ausklammern (Einheit 2), binomische Formeln (Einheit 4), Quadratwurzeln}
\setcounter{aufgabe}{39}

\begin{aufgabe}{Ich erkenne, ob ein Term ein Quadrat ist.}
\anweisung{Ist der Term ein Quadrat? Kreuze an.}
\begin{teilezwei}
\tz $y^2$ \quad \janein & \tz $64$ \quad \janein \\
\tz $25x^2$ \quad \janein & \tz $12x$ \quad \janein \\
\tz $18$ \quad \janein & \tz $3x^2$ \quad \janein \\
\tz $100y^2$ \quad \janein & \\
\end{teilezwei}
\end{aufgabe}

\begin{aufgabe}{Ich erkenne die beiden Quadrate und das Mittelglied.}
\anweisung{Unterstreiche die beiden Quadrate und kreise das Mittelglied ein. Rechne nichts aus.}
\begin{teilezwei}
\tz \rule{0pt}{9mm}{\large $y^2 + 16y + 64$} & \tz \rule{0pt}{9mm}{\large $36 + 12x + x^2$} \\
\tz \rule{0pt}{9mm}{\large $9x^2 - 24x + 16$} & \tz \rule{0pt}{9mm}{\large $49 - 14y + y^2$} \\
\tz \rule{0pt}{9mm}{\large $81x^2 + 18x + 1$} & \\
\end{teilezwei}
\end{aufgabe}

\verfahren{Faktorisieren mit den Formeln}

\begin{aufgabe}{Ich kann einen Term mit einer binomischen Formel faktorisieren.}
\anweisung{Faktorisiere.}
\rechenplatz{2}
\begin{gleichungsraster}[1]
\gl{x^2 - 16} & \gl{x^2 - 36} \\
\gl{y^2 - 100} & \gl{x^2 - 1} \\
\gl{x^2 + 12x + 36} & \gl{x^2 - 16x + 64} \\
\gl{9x^2 + 30x + 25} & \gl[2]{2x^2 - 72} \\
\anweisung{Faktorisiere und mache die Probe durch Ausmultiplizieren.}
\gl[2]{x^2 - 20x + 100} & \gl[3]{5x^3 - 40x^2 + 80x} \\
\end{gleichungsraster}
\end{aufgabe}

\verfahren{Gleichungen durch Faktorisieren lösen}

\begin{aufgabe}{Ich kann eine Gleichung lösen, indem ich faktorisiere.}
\anweisung{Faktorisiere und löse die Gleichung.}
\rechenplatz{3}
\begin{gleichungsraster}[3]
\gl{x^2 - 81 = 0} & \gl{x^2 - 144 = 0} \\
\gl{x^2 - 18x + 81 = 0} & \gl{3x^2 - 48 = 0} \\
\end{gleichungsraster}
\end{aufgabe}

\verfahren{Quadratische Ergänzung}

\begin{aufgabe}{Ich kann einen Term als Quadrat einer Klammer plus eine Zahl schreiben.}
\anweisung{Schreibe den Term als Quadrat einer Klammer plus oder minus eine Zahl.}
\rechenplatz{3}
\begin{gleichungsraster}[2]
\gl{x^2 + 4x + 7} & \gl{x^2 + 8x + 20} \\
\gl{x^2 - 12x + 40} & \gl{x^2 + 16x + 50} \\
\end{gleichungsraster}
\end{aufgabe}

\begin{aufgabe}{Ich finde den Fehler beim Faktorisieren.}
\anweisung{Finde den Fehler, benenne ihn und korrigiere.}
\begin{teile}
\teil Jana faktorisiert: \rechnung{x^2 + 64 &= (x + 8) \cdot (x - 8)}
\teil Ali faktorisiert: \rechnung{x^2 + 9x + 16 &= (x + 4)^2}
\teil Tim faktorisiert: \rechnung{2x^2 - 32 &= 2 \cdot (x^2 - 32)}
\end{teile}
\end{aufgabe}

\begin{aufgabe}{Ich kann entscheiden, ob sich ein Term mit einer binomischen Formel faktorisieren lässt.}
\anweisung{Lässt sich der Term mit einer binomischen Formel faktorisieren? Kreuze an, ohne zu rechnen.}
\begin{teilezwei}
\tz $x^2 + 22x + 121$ \quad \janein & \tz $x^2 + 36$ \quad \janein \\
\tz $y^2 - 169$ \quad \janein & \tz $x^2 + 5x + 4$ \quad \janein \\
\tz $4x^2 - 20x + 25$ \quad \janein & \tz $x^2 - 8x - 16$ \quad \janein \\
\end{teilezwei}
\end{aufgabe}

\begin{aufgabe}{Ich kann begründen, wann ein Term keine binomische Formel ist.}
\begin{teile}
\teil Begründe, warum sich $x^2 - 100$ mit einer binomischen Formel faktorisieren lässt, $x^2 + 100$ aber nicht.
\teil $x^2$ und $36$ sind Quadrate. Begründe, warum $x^2 + 10x + 36$ trotzdem keine binomische Formel ist.
\end{teile}
\end{aufgabe}

\begin{aufgabe}{Ich kann eine Sachaufgabe durch Faktorisieren lösen.}
Eine quadratische Terrasse hat den Flächeninhalt $(x^2 + 18x + 81)$\,m$^2$.

\rechteck[seiten={?,?,,},punkte=]{3}{3}

\begin{teile}
\teil Wie lang ist eine Seite der Terrasse? Faktorisiere dazu.
\teil Gib einen Term für den Umfang der Terrasse an.
\teil Für die Terrasse gilt $x = 3$. Wie lang ist eine Seite, und wie groß ist die Terrasse?
\end{teile}
\end{aufgabe}
